\subsection{永田定理}

定理 1（Tate）．——设 A 是诺特整闭环，a 是 A 的一个元素。假设理想 aA 是素理想，环 A/aA 是日式环，且 A 关于 aA-进拓扑完备。则环 A 是日式环。

\begin{enumerate}
\def\labelenumi{\alph{enumi})}
\tightlist
\item
  令 K 为 A 的分式域。当 K 的特征为 0 时断言是平凡的（第 1 号，命题 1 的推论），故可假设 K 的特征为 \(p > 0\)。也可假设 \(a \neq 0\)。
\end{enumerate}

令 L 为 K 的有限根式扩张，q 为 p 的一个幂，使 \(L \subset K^{1/q}\)。置 \(x = a^{1/q}\)，\(M = L(x)\)。根据第 1 号的命题 1，只需证明 A 在 M 中的整闭包 B 是有限 A-代数。

\begin{enumerate}
\def\labelenumi{\alph{enumi})}
\setcounter{enumi}{1}
\item
  先证明理想 \(xB\) 是 B 中位于 \(aA\) 之上的唯一素理想。事实上，B 中至少存在一个位于 \(aA\) 之上的素理想（V，§ 2，第 1 号，定理 1）。令 \(\mathfrak{q}\) 为其中一个理想。有 \(x^{q} = a \in \mathfrak{q}\)，从而 \(xB \subset \mathfrak{q}\)，因为 \(\mathfrak{q}\) 是素理想。反之，取 \(y\) 为 \(\mathfrak{q}\) 中的元素；K 中的元素 \(y^{q}\) 在 A 上整，由于 A 整闭，故属于 A。由于 \(\mathfrak{q} \cap A = aA\)，存在 A 中的元素 \(\alpha\)，使 \(y^{q} = a\alpha = x^{q}\alpha\)。因此，M 中的元素 \(y/x\) 在 A 上整，故属于 B；于是有 \(y \in xB\)，从而 \(\mathfrak{q} = xB\)，断言得证。
\item
  由此，环 \(B_{xB}\) 是环 \(A_{aA}\) 在 M 中的整闭包（V，§ 1，第 5 号，命题 16 及 § 2，第 1 号，命题 2）。根据 VI，§ 3，第 6 号，命题 9，\(A_{aA}\) 是离散赋值环；于是由 Krull-Akizuki 定理（VII，§ 2，第 5 号，命题 5）可推出，域 \(\kappa(xB)\) 是 \(\kappa(aA)\) 的有限扩张，且 \(B_{xB}\) 是诺特环。
\item
  环 B/xB 在日式环 A/aA 上整，且其分式域是后者分式域的有限扩张。因此，B/xB 是有限生成的 (A/aA)-模。对于每个整数 \(i \geqslant 0\)，模 \(x^{i}B/x^{i+1}B\) 也具有同样性质。因此，(A/aA)-模 B/aB 有一个长度为 q 的合成列，其商都是有限生成的 (A/aA)-模，因而 B/aB 本身是有限生成的 (A/aA)-模。
\item
  给环 A 赋予 (aA)-进滤子，给环 B 赋予 (aB)-进滤子。根据假设 A 是完备的；由于 \(B_{xB}\) 是整诺特环，其 \(aB_{xB}\)-进滤子在 \(B_{xB}\) 上是分离的（III，§ 3，第 2 号，命题 5 的推论），所以有 \(\bigcap a^{n}B \subset \bigcap a^{n}B_{xB} = \{0\}\)，且 B 的 \(aB\)-进滤子是分离的；gr(A)-模 gr(B) 由 gr₀(B) 生成，故根据 d) 它是有限生成的。于是根据 III，§ 2，第 9 号，命题 12 的推论 1，B 是有限生成的 A-模，证明完毕。
\end{enumerate}

推论．——设 R 是诺特整环，n 是一个整数。若 R 是日式环，则环 \(R[[T_{1}, \ldots, T_{n}]]\) 是日式环。

由归纳法可设 \(n = 1\)。令 S 表示 R 的整闭包；若 R 是日式环，则 S 是 R 上的有限代数，因而是日式环（第 1 号，命题 2）。环 \(S[[T]]\) 是诺特环且整闭（V，§ 1，第 4 号，命题 14）；将定理 1 应用于 A = \(S[[T]]\) 和 \(a = T\)，得到 \(S[[T]]\) 是日式环。因此 \(R[[T]]\) 是日式环（第 1 号，命题 2）。

定理 2（Nagata）．——每个完备局部诺特整环 A 都是日式环。

根据第 2 节第 5 号的定理 3 以及第 3 节第 3 号的定理 2，存在整数 \(n \geqslant 0\)、一个或者是域或者是其分式域特征为 0 的离散赋值环的环 R，以及 A 的一个子环 B；B 同构于 \(R[[T_{1}, \ldots, T_{n}]]\)，且 A 是有限 B-代数。于是 R 是日式环（第 1 号，命题 1 的例子和推论），从而 B 是日式环（定理 1 的推论），A 也是日式环（第 1 号，命题 2）。

推论．——设 A 是一个诺特半局部环，其完备化是约化的。则 A 在其全分式环 R 中的整闭包 A' 是有限 A-代数。

先设 A 是局部且完备的，并设 \(p_{1}, \ldots, p_{n}\) 是 A 的互异极小素理想；对 \(i = 1, \ldots, n\)，记 \(K_{i}\) 为 \(A/p_{i}\) 的分式域，记 \(A_{i}'\) 为 \(A/p_{i}\) 的整闭包。由于 A 是约化的，R 是环 \(K_{i}\) 的乘积，而 \(A'\) 是环 \(A_{i}'\) 的乘积（V，§ 1，第 2 号，命题 9 的推论 1）。由于各局部环 \(A/p_{i}\) 都是整环且完备，它们是日式环（定理 2），所以每个 \(A_{i}'\) 都是有限 A-代数，\(A'\) 也是有限 A-代数。

若 A 是半局部且完备的，则它同构于有限个完备局部环的乘积（III，§ 2，第 13 号，命题 19 的推论），由前述结果立即得到结论。

转到一般情形。注意，A 的完备化 \(\hat{\mathbf{A}}\) 是一个半局部、完备的诺特环，并且忠实平坦于 A（III，同上引文，§ 3，第 4 号，命题 8 的推论以及 § 3，第 5 号，命题 9）。令 S 为 A 的非零因子集合；有 R = S \(^{-1}\) A。由于 \(\hat{\mathbf{A}}\) 平坦于 A，S 的元素在 \(\hat{\mathbf{A}}\) 中都是非零因子，并且 S \(^{-1}\) \(\hat{\mathbf{A}}\) 可识别为 \(\hat{\mathbf{A}}\) 的全分式环 T 的一个子环。仍由于 \(\hat{\mathbf{A}}\) 平坦于 A，环 A′ ⊗ \(_{\mathbf{A}}\) \(\hat{\mathbf{A}}\) 可识别为 R ⊗ \(_{\mathbf{A}}\) \(\hat{\mathbf{A}}\) = S \(^{-1}\) \(\hat{\mathbf{A}}\) 的一个子环，因而也可识别为 T 中整于 \(\hat{\mathbf{A}}\) 的一个子环。根据证明的第一部分，A′ ⊗ \(_{\mathbf{A}}\) \(\hat{\mathbf{A}}\) 因而是一个 \(\hat{\mathbf{A}}\) -有限型模；所以 A′ 是有限生成的 A-模（I，§ 3，第 6 号，命题 11）。

回顾（A，V，第 114 页，定义 1），域 K 上的代数 E 称为可分的，是指对 K 的每个扩张 L，环 L \(\otimes_{\mathbf{K}}\) E 都是约化的；只需对 K 的每个有限扩张成立即可。下面的命题推广了定理 2：

命题 3．——设 A 是诺特半局部整环，K 是其分式域。若 K-代数 \(K \otimes_{A} \hat{\mathbf{A}}\) 是可分的，则环 A 是日式环。

设 L 是 K 的有限扩张，B 是 A 在 L 中的整闭包。取 B 的有限子集 F，使 \(L = K[F]\)（V，§ 1，第 5 号，命题 16 的推论 2）；令 C 为 F 生成的（有限）A-代数。由于 L 是 C 的分式域，环 B 是 C 的整闭包（V，§ 1，第 1 号，命题 6），所以只需证明 B 是有限 C-代数。现在，C 是半局部诺特环（IV，§ 2，第 5 号，命题 9 的推论 3）；其完备化可识别为 \(C \otimes_{A} \hat{A}\)（III，§ 3，第 4 号，定理 3 (ii)），因而也可识别为约化环 \(L \otimes_{A} \hat{A} = L \otimes_{K} (K \otimes_{A} \hat{A})\) 的一个子环，于是它是约化的。因此命题 3 由定理 2 的推论得到。

